\section{Dataset Details and Prompt Templates}
\label{sec:appendix}

\subsection{Dataset Details}

POPE (Polling-based Object Probing Evaluation)
is a benchmark built on the COCO dataset
\citep{linMicrosoftCOCOCommon2015}
that converts object hallucination evaluation
into a binary existence probing task
\citep{liEvaluatingObjectHallucination2023}.
The paradigm was originally proposed
by \citet{rohrbachObjectHallucinationImage2018}
for image captioning hallucination
and later systematized by \citet{liEvaluatingObjectHallucination2023}.
Given an image, the benchmark constructs binary questions
asking whether a candidate object exists.
Three negative sampling strategies create
an increasing difficulty gradient:
Random (random negatives, easiest),
Popular (high-frequency co-occurring but absent objects),
and Adversarial (semantically similar but absent objects, hardest).
This experiment uses all three splits,
1{,}000 items each, totaling 3{,}000 per model.

MathVista is a comprehensive benchmark
for evaluating mathematical reasoning
of foundation models in visual contexts
\citep{luMathVistaEvaluatingMathematical2024}.
Questions span chart interpretation, geometry, numerical computation,
and function analysis,
in both multiple-choice and free-response formats.
The task requires simultaneous visual perception,
mathematical knowledge, and multi-step reasoning,
making it well suited to trigger and analyze all three hallucination types.
This experiment uses the testmini subset
(approximately 1{,}000 samples).

VQA-RAD is a standard medical VQA benchmark
\citep{lauDatasetClinicallyGenerated2018}
containing 2{,}244 clinically generated question-answer pairs
over 314 radiology images (X-ray, CT, MRI).
Answers are either closed-ended (yes/no, approximately 56\%)
or open-ended (free-text diagnosis, approximately 44\%).
This experiment uses the test split (451 samples).

%% ===================================================================
\subsection{POPE Response Generation Prompt}

The model prompt for POPE concatenates the question
with a fixed format instruction.
\Cref{fig:pope_prompt_example} shows a concrete example.

\begin{figure}[htbp]
  \begin{minipage}{0.30\textwidth}
    \centering
    \includegraphics[width=\linewidth]{pope_chair_sample.jpg}
  \end{minipage}\hfill
  \begin{minipage}{0.65\textwidth}
    Prompt sent to the model:
\begin{lstlisting}
Is there a chair in the image?
Please answer YES or NO without
an explanation.
\end{lstlisting}
  \end{minipage}
  \caption{POPE prompt example
    (COCO image with a chair-related probing question).}
  \label{fig:pope_prompt_example}
\end{figure}

%% ===================================================================
\subsection{MathVista Response Generation Prompt}

MathVista prompts are dynamically constructed from sample metadata.
The format instruction varies by question type:
``Answer with the option letter only.'' for multiple-choice,
``Answer with an integer.'' for integer answers,
``Answer with a number rounded to $k$ decimal place(s).'' for float answers,
and ``Keep your answer concise.'' as the default.
If choices are present, they are appended on a separate line.

In Direct mode, the prompt structure is:

\begin{lstlisting}
Question: {query}

{format_instruction}
Choices: {choices}
\end{lstlisting}

In CoT mode, a step-by-step reasoning preamble is prepended:

\begin{lstlisting}
Please generate a step-by-step answer,
base your reasoning strictly on what is
visible in the image, and avoid speculating
about details that are unclear or not
present. End with the final answer on a
separate line starting with 'Final answer:'.

Question: {query}

{format_instruction}
Choices: {choices}
\end{lstlisting}

\Cref{fig:mathvista_prompt_example} shows
a concrete CoT-mode example using a geometry problem
from MathVista testmini (pid=5, question\_type=multi\_choice).

\begin{figure}[htbp]
  \begin{minipage}{0.30\textwidth}
    \centering
    \includegraphics[width=\linewidth]{mathvista_pid5.jpg}
  \end{minipage}\hfill
  \begin{minipage}{0.65\textwidth}
    CoT prompt sent to the model:
\begin{lstlisting}
Please generate a step-by-step answer,
base your reasoning strictly on what is
visible in the image, and avoid speculating
about details that are unclear or not
present. End with the final answer on a
separate line starting with 'Final answer:'.

Question: Find $m\angle H$

Answer with the option letter only.
Choices: 97, 102, 107, 122
\end{lstlisting}
  \end{minipage}
  \caption{MathVista CoT prompt example
    (geometry problem, pid=5).}
  \label{fig:mathvista_prompt_example}
\end{figure}

%% ===================================================================
\subsection{VQA-RAD Response Generation Prompt}

The VQA-RAD prompt adopts an expert role preamble
following the strategy of
\citet{panditHALTMedVQAHallucinationAware2025}.
\Cref{fig:vqarad_prompt_example} shows a concrete example.

\begin{figure}[htbp]
  \begin{minipage}{0.30\textwidth}
    \centering
    \includegraphics[width=\linewidth]{vqarad_sample.jpg}
  \end{minipage}\hfill
  \begin{minipage}{0.65\textwidth}
    Prompt sent to the model:
\begin{lstlisting}
You are an expert radiologist.

Is there evidence of an aortic aneurysm?
Please provide a detailed answer based on
what you observe in the medical image.
\end{lstlisting}
  \end{minipage}
  \caption{VQA-RAD prompt example (chest X-ray).}
  \label{fig:vqarad_prompt_example}
\end{figure}

%% ===================================================================
\subsection{MLLM Judge System Prompt}

The judge system prompt consists of three parts:
a rating scale definition,
few-shot examples with an important reminder,
and a hallucination type classification with output format.
The complete text is reproduced below.

\begin{lstlisting}
You are an impartial evaluator for multimodal
language models. Your task is to rate a model's
response to a visual question on a 0-6 scale,
jointly considering informativeness and
hallucination.

## Rating Scale (0-6)
- 6: Very informative, with good
     analysis/reasoning, NO hallucination
- 5: Very informative, NO hallucination
- 4: Somewhat informative, NO hallucination
- 3: Not informative (e.g., "I don't know"),
     but NO hallucination
- 2: Very informative, WITH hallucination
- 1: Somewhat informative, WITH hallucination
- 0: Not informative, WITH hallucination

## Hallucination Definition
Hallucination refers to situations where the
response includes information not present or
implied in the image, or makes claims
inconsistent with the visual content.
A response that honestly admits inability to
answer (e.g., "I cannot determine") has NO
hallucination (score 3).

## Few-Shot Examples

Example 1 - Hallucination (Score 2):
Image shows two cats. Question: "How many cats
are in the image?" Model response: "There is
one cat in the image." The response is
informative but miscounts a visible object.
Score: 2.

Example 2 - No Hallucination (Score 5):
Image shows a panda sitting on a tree branch
with green leaves. Question: "What is the
panda doing?" Model response: "The panda is
sitting on a tree branch, surrounded by
leaves." The response is accurate and makes no
false claims. Score: 5.

Example 3 - Hallucination (Score 2):
Image shows a dog holding an umbrella, with no
water in sight. Question: "Describe the scene."
Model response: "A dog holding an umbrella,
positioned near the water." The response
fabricates "water" which is not in the image.
Score: 2.

Example 4 - No Hallucination (Score 3):
Image shows a complex diagram or chart.
Question: "Explain the economic trend." Model
response: "I cannot determine the specific
trend from this image. The chart appears to
show some data over time." The response
honestly admits uncertainty without making
false claims. Score: 3.

Example 5 - No Hallucination (Score 6):
Image shows industrial smokestacks. Question:
"What environmental concerns does this image
raise?" Model response provides a detailed
analysis of potential emissions, air quality
impacts, and relevant environmental
regulations, all grounded in visible evidence.
Score: 6.

## Important Reminder
The standard human-generated answer may only
contain factual information without detailed
analysis, and may not be fully comprehensive.
A model response that provides detailed
reasoning or analysis grounded in the visual
evidence should be encouraged, not penalized.
Only penalize claims that are clearly
inconsistent with or absent from the image.

## Hallucination Types (only when score < 3)
- faithfulness: Visual inconsistency -
  describes objects/attributes/relations not
  in the image
- factuality: Contradiction with established
  world knowledge
- logical: Reasoning error - conclusion does
  not follow from evidence or
  self-contradiction

## Output Format
Reply in JSON only:
{"score": 0-6, "has_hallucination": true/false,
 "hallucination_type": "faithfulness" |
   "factuality" | "logical" | "none",
 "reason": "Brief explanation in English"}
\end{lstlisting}

%% ===================================================================
\subsection{MLLM Judge User Prompt}

For each evaluation, the judge receives a user prompt
constructed from the task description, question,
model response, and ground truth.
The task description varies by dataset:

\begin{lstlisting}
# POPE
The question is a Yes/No object existence
query.

# MathVista
The question requires mathematical visual
reasoning (chart, diagram, or geometry).

# VQA-RAD
The question is about medical visual question
answering based on a radiology image (X-ray,
CT, or MRI). Carefully check whether the
model's response is consistent with the
visual content of the image and with
established medical and anatomical knowledge.
\end{lstlisting}

The user prompt template is:

\begin{lstlisting}
{task_desc}

Question: {question}
Model Response: {model_response}
Ground Truth Answer: {ground_truth}

Please rate the model's response on the
0-6 scale.
\end{lstlisting}

\Cref{fig:judge_prompt_example} shows
a concrete example using the MathVista geometry problem (pid=5),
where GPT-5.4-mini produced a correct CoT response (score 6).

\begin{figure}[htbp]
  \begin{minipage}{0.30\textwidth}
    \centering
    \includegraphics[width=\linewidth]{mathvista_pid5.jpg}
  \end{minipage}\hfill
  \begin{minipage}{0.65\textwidth}
    User prompt sent to the judge:
\begin{lstlisting}
The question requires mathematical visual
reasoning (chart, diagram, or geometry).

Question: Find $m\angle H$
Choices: 97, 102, 107, 122
Model Response: The figure shows a
pentagon with interior angles labeled:
angle E = x, angle F = (x+20),
angle G = (x+5), angle H = (x-5),
angle J = (x+10). For a pentagon, the
sum of interior angles is 540. So
x + (x+20) + (x+5) + (x-5) + (x+10)
= 540, 5x + 30 = 540, x = 102.
Then m angle H = x - 5 = 97.
Final answer: A
Ground Truth Answer: 97

Please rate the model's response on the
0-6 scale.
\end{lstlisting}
  \end{minipage}
  \caption{MLLM Judge user prompt example (MathVista pid=5).
    The judge returns
    \texttt{\{"score": 6, "has\_hallucination": false,
    "hallucination\_type": "none", "reason": "..."\}}.}
  \label{fig:judge_prompt_example}
\end{figure}
